% Representational Incompleteness — LaTeX source (Lean 4 / Mathlib companion library)

\documentclass[11pt]{article}
\usepackage[margin=1in]{geometry}
\usepackage{amsmath,amssymb,amsthm}
\usepackage{hyperref}
\usepackage{enumitem}
\usepackage{array}
\usepackage{booktabs}
\usepackage{xcolor}
\usepackage{stmaryrd}

% Allow line breaks in long Lean identifiers (\_, namespaces, CamelCase).
\ExplSyntaxOn
\tl_new:N \l__rr_leanref_tl
\NewDocumentCommand \leanref { m }
  {
    \tl_set:Nn \l__rr_leanref_tl { #1 }
    \regex_replace_all:nnN { \c{_} } { \c{_} \c{allowbreak} } \l__rr_leanref_tl
    \regex_replace_all:nnN { \. } { . \c{allowbreak} } \l__rr_leanref_tl
    % CamelCase / digit-boundary splits (Lean module and declaration names).
    \regex_replace_all:nnN { ([a-z])([A-Z]) } { \1 \c{allowbreak} \2 } \l__rr_leanref_tl
    \regex_replace_all:nnN { ([0-9])([A-Z]) } { \1 \c{allowbreak} \2 } \l__rr_leanref_tl
    \texttt{\small \tl_use:N \l__rr_leanref_tl }
  }
\ExplSyntaxOff

\newtheorem{theorem}{Theorem}[section]
\newtheorem{lemma}[theorem]{Lemma}
\newtheorem{proposition}[theorem]{Proposition}
\newtheorem{corollary}[theorem]{Corollary}
\newtheorem{definition}[theorem]{Definition}
\newtheorem{remark}[theorem]{Remark}

\newcommand{\novel}{\textcolor{blue!70!black}{[Novel]}}
\newcommand{\known}{\textcolor{gray}{[Known math]}}
% Break `theorem`/`definition` head away from long \leanref tags (avoids overfull italic headers).
\newcommand{\TheoremLeanTag}[1]{%
  \noindent\textnormal{(\leanref{#1})}\par\noindent}
% End `amsthm` head paragraph so the declaration tag is not merged into the header line.
\newcommand{\AfterTheoremHead}{\mbox{}\par}

\title{Representational Incompleteness:\\
  No Self-Modeling System Can Capture Its Own Diagonal\\[0.3em]
  {\large A Machine-Checked Barrier with Six Closed Escape Routes}}
\author{Nova Spivack\\
Independent Researcher}
\date{2026}

\begin{document}
\maketitle

% ============================================================================
% ABSTRACT
% ============================================================================
\begin{abstract}
\noindent\textbf{Summit theorem (\leanref{representational\_incompleteness}).}
Any system that models itself---any $s : A \to A \to B$ whatsoever---is
\emph{necessarily incomplete} with respect to its own diagonal behavior:
the function $a \mapsto f(s\,a\,a)$ is never equal to any row $s\,a_0$
of the self-model, for every fixed-point-free $f : B \to B$ and every
$a_0 \in A$.  No surjectivity, totality, or cardinality assumption on
$s$ is required.  If $s$ is claimed to be total/surjective, it simply
cannot exist.

This is \textbf{representational incompleteness}: the analogue of
G\"odel incompleteness, not for formal theories and provability, but
for any system that models itself via a parametric family of functions.
The blind spot is structural---not a resource limitation---and
persists for every codomain admitting a fixed-point-free endomorphism
(including $\mathbb{N}$ with successor and $\mathrm{Bool}$ with
negation).

\medskip
\noindent\textbf{Topological reinforcement.}
A second, independent strand proves that the barrier is not merely
diagonal but \emph{topological}: a general lemma states that any two
spaces whose intrinsic planar-chart boundaries are subtype-incompatible
are not homeomorphic.  The obstacle is \textbf{wrong type}, not
\textbf{inexhaustible queue}---a finite-data invariant that no amount
of computation, scaling, or architectural extension can change.  The
lemma applies to an \emph{open-ended class} of target spaces (M\"obius
strip, Klein bottle, projective plane, or any non-orientable surface
with incompatible boundary); the closed cylinder and M\"obius strip are
the minimal fully discharged compact witness.

\medskip
\noindent\textbf{No escape.}
Six adversarial retreat routes are formally closed: infinite regress,
finite-order iterate collapse, cross-object readout (both sides),
partial self-models, decidable self-reference, and model substitution.
Each is a proved theorem in Lean~4, checked against Mathlib, with zero
\texttt{sorry} and no project-specific axioms.

\medskip
\noindent\textbf{Consequence.}
In every neighboring debate---consciousness, symbolic AI, semantic
grounding, measurement chains, LLM scaling---the six closed escape routes
leave only one admissible shape for solutions: \emph{structural change}
to the topological type of the substrate, not quantitative extension
within a fixed type.  Scaling, iterating, partial modeling, cross-object
readout, finite-depth loops, and decidable restriction are all formally
inadmissible.  Self-modeling systems are provably incomplete; the
formalism makes this precise and machine-checkable.
\end{abstract}

% ============================================================================
% 1. INTRODUCTION
% ============================================================================
\section{Introduction}
\label{sec:intro}

\subsection{The Problem and the Summit Theorem}
When a system models itself---by iterating a fixed construction, by
encoding its own states, by interpreting its own symbols---a familiar worry
arises: infinite regress, or diagonal limitations, or both.  But those
concerns are usually framed \emph{computationally}: the regress does not
terminate, the self-interpreter cannot be total, the G\"odel sentence is
undecidable.  The implicit hope is that with enough resources---more
compute, richer architectures, deeper self-models---the obstacle might
shrink or be routed around.

This paper proves it cannot.  The headline result is:

\begin{theorem}[\novel\ Representational Incompleteness]
\label{thm:summit}
\AfterTheoremHead
\TheoremLeanTag{representational\_incompleteness}
Let $B$ be a type with a fixed-point-free endomorphism $f : B \to B$.
For \textbf{any} $s : A \to A \to B$ (no surjectivity, totality, or
cardinality assumption):
\begin{enumerate}[label=(\alph*)]
\item \textbf{Diagonal exclusion:}
  for every $a_0 \in A$,
  $\;s\,a_0 \neq \lambda a.\,f(s\,a\,a)$.
\item \textbf{Total coverage impossible:}
  $\neg\!\left(\forall g : A \to B,\;\exists a,\;s\,a = g\right)$.
\end{enumerate}
\end{theorem}

\noindent
Every self-modeling system is \emph{necessarily incomplete}: its own
diagonal behavior is precisely the information it cannot represent.  If
$s$ is claimed to be total and surjective, it simply cannot exist.  This
is \textbf{representational incompleteness}---the analogue of G\"odel
incompleteness, not for formal theories and provability, but for any
system that models itself via a parametric family of functions.  The
blind spot is structural, not a resource limitation, and persists for
$B = \mathbb{N}$ (successor), $B = \mathrm{Bool}$ (negation), and every
other codomain admitting a fixed-point-free endomorphism.

A second, independent strand reinforces this with a
\textbf{topological} barrier: a general lemma showing that
\emph{any} two spaces whose intrinsic planar-chart boundaries
are subtype-incompatible are not homeomorphic---an open-ended
class of obstructions that no amount of iteration, scaling, or
continuous deformation within a fixed topological type can
remove.  And we close, one by one, every standard adversarial
route by which one might hope to circumvent either barrier.

The technical development is in Lean~4, checked against Mathlib, with
zero \texttt{sorry} and no project-specific axioms;
\S\ref{sec:instantiations}--\S\ref{sec:admissibility} relate the
results to consciousness, symbolic grounding, diagonal arguments,
AI scaling, and observer--apparatus topologies.

\subsection{Summary of Results}

The core results are \textbf{representational incompleteness}
(Theorem~\ref{thm:summit}) and a \textbf{universal topological
obstruction}---reinforced by a proof that the six most natural
adversarial attempts to evade them all fail.

\paragraph{The topological obstruction.}
A general lemma (Theorem~\ref{thm:generic}) states: if two spaces admit
subsets characterized as exactly the non-\texttt{ChartableR2} points, and
if those boundary \emph{subspaces} are not homeomorphic, then the ambient
spaces are not homeomorphic.  This is a \emph{finite-data invariant}: two
spaces are either homeomorphic or not, regardless of how many modeling
steps one enumerates.  The lemma defines an \emph{open-ended class} of
obstructed pairs---any two spaces meeting the three hypotheses, regardless
of genus, orientability, or dimension---rather than a single hard-wired
example.  The closed cylinder and M\"obius strip serve as the
\textbf{minimal-case witness} (\textbf{M-FINAL}); Klein bottles,
projective planes, higher-genus non-orientable surfaces, etc., would
equally suffice.  Any future domain whose state space matches the same
boundary pattern imports the same conclusion with no new proof.

\paragraph{Why this is not ``just'' infinite regress.}
The categorical strand (\S\ref{sec:regress}) does yield infinitely many
distinct meta-levels under the bundled hypothesis
\texttt{iter\_injective}, with Lawvere's wall and
\texttt{OntologicalSlot} preventing collapse of the endo/obj distinction.
But infinity is not the deepest point.  The topological obstruction is
independent: it says the barrier is \textbf{wrong type}, not
\textbf{inexhaustible queue}.  No amount of compute, no deeper
self-model, no richer architecture changes the boundary pattern of a
topological type.

\paragraph{No escape.}
Six adversarial retreat routes are formally closed
(\S\ref{sec:noescape}):
\begin{enumerate}[label=(\roman*),leftmargin=2em,itemsep=0.15ex,topsep=0.5ex,parsep=0pt]
\item Infinite injective tower $\Rightarrow$ full regress + Lawvere wall.
\item Finite-order iterate collapse $\Rightarrow$ bounded depth, no
  injective tower.
\item Cross-object readout via section--retraction $\Rightarrow$
  idempotent on host, identity on source; no injective tower on
  \emph{either} side.
\item Partial (non-surjective) self-model $\Rightarrow$ diagonal blind
  spot---the very information needed for self-capture is the information
  it cannot represent.
\item Decidable / $\mathrm{Bool}$-valued self-reference $\Rightarrow$
  Lawvere obstruction persists in the finite, two-valued setting.
\item Different topological model $\Rightarrow$ the general lemma
  applies to \emph{any} pair with incompatible \texttt{ChartableR2}
  boundaries, not just cylinder/M\"obius.
\end{enumerate}
Together these leave an adversary with a clean dichotomy: either an
infinite injective regress with the Lawvere wall intact, or a finite
loop too shallow for full self-capture---and in either case the
topological barrier stands.

\paragraph{Bridges and conditional readings.}
Cross-domain structures
(\leanref{SymbolSystem},
\leanref{LawvereRegressBridge},
bundled \leanref{ChartableR2BoundaryModel} instances,
\leanref{QuantumObserverChainHypothesis})
make the obstruction's \emph{hypothesis shape} explicit for grounding,
diagonal, consciousness~(under premise~P2), and observer--readout
readings.  Because the topological obstruction is a general lemma over
a class of space-pairs, each bridge inherits the full generality
automatically.  An extended master theorem adjoins all three layers
(diagonal, topological, categorical) in one checkable bundle.

\paragraph{Utility.}
For any application the workflow is:
(1)~encode the system's state space in sufficient topological detail to
verify the boundary biconditionals,
(2)~compare boundary subspaces up to homeomorphism,
(3)~invoke the general lemma or a bundled model.
If the premises hold, the obstruction applies; if they fail, the
obstacle was never licensed.  That conditional structure is what keeps
the mathematics honest while delivering a \emph{portable} gate.

\subsection{Contributions and Novelty Claims}

Individual ingredients are classical: Lawvere's diagonal (1969), textbook
topology of surfaces, Cantorian regress patterns.
The contribution is a \textbf{unified, adversarially checked, machine-verified}
landscape in which representational incompleteness, topological obstruction,
and six closed escape routes coexist in one artifact, so that ``compute
harder'' is never confusable with ``change homeomorphism type.''  Specific
novel elements:

\begin{enumerate}[label=(\arabic*),itemsep=0.4ex]
\item \textbf{\texttt{OntologicalSlot}} --- inductive tagging of morphism
  vs.\ global-element slots; makes the Lawvere wall explicit inside the
  type discipline. [\novel]
\item \textbf{\texttt{ChartableR2} spine + seam--equator charts} ---
  intrinsic $\mathbb{R}^2$ chartability as a definable boundary sensor;
  constructive charts through the M\"obius quotient seam including the
  equator via piecewise unfolding.  The M\"obius strip serves as the
  minimal proof witness; the predicate applies uniformly to all
  surfaces. [\novel\ technique]
\item \textbf{General boundary-contrast obstruction} ---
  a single reusable, machine-checked lemma defining an
  \emph{open-ended class} of obstructed pairs: any two spaces
  with incompatible \leanref{ChartableR2}-co-characterized
  boundaries are not homeomorphic---regardless of genus,
  orientability, or dimension.  To our knowledge, \emph{not}
  standard in the philosophy-of-mind or AI-safety
  literatures. [\novel]
\item \textbf{Cross-domain bridge structures} ---
  \texttt{SymbolSystem}, \texttt{LawvereRegressBridge},
  bundled \texttt{ChartableR2BoundaryModel} data,
  \texttt{QuantumObserverChainHypothesis}: these make the obstruction's
  hypothesis shape explicit per domain. [\novel]
\item \textbf{No-escape horn closures} --- six formally blocked
  adversarial escape routes (iterate collision, section--retraction on
  both sides, partial-model diagonal exclusion, decidable Lawvere,
  wall persistence at every depth).  The integrated closure package is,
  to our knowledge, new. [\novel]
\end{enumerate}

\subsection{Paper Organization}
The paper is structured around three layers that build toward a single
conclusion.

\textbf{Layer~1: Diagonal \& categorical}
(\S\ref{sec:regress})---representational systems, regress, Lawvere wall,
\texttt{OntologicalSlot}, iterate dichotomy, core master theorem.

\textbf{Layer~2: Topological obstruction}
(\S\S\ref{sec:models}--\ref{sec:mfinal})---minimal-case witnesses
(cylinder, M\"obius strip), \texttt{ChartableR2} boundary sensor, general
obstruction lemma, seam charts, M-FINAL.

\textbf{Layer~3: Integration \& closure}
(\S\ref{sec:extended})---extended master theorem, instantiations
(consciousness, grounding, diagonal arguments, QM, LLM scaling),
horn-by-horn no-escape analysis~(\S\ref{sec:noescape}),
admissibility~(\S\ref{sec:admissibility}).

\S\S\ref{sec:limits}--\ref{sec:conclusion}: limitations, related work,
conclusion.  Appendices: Lean module map, key theorem index, reproduction.

% ============================================================================
% 2. CATEGORICAL STRAND: THE REPRESENTATIONAL REGRESS
% ============================================================================
\section{The Representational Regress}
\label{sec:regress}

This section builds the categorical machinery that feeds into the summit
theorem: the endomorphism tower, the Lawvere wall, and the
\texttt{OntologicalSlot} distinction.  Together they establish that
self-representation via iterated endomorphisms either produces an
infinite regress with the endo/obj wall intact, or collapses to a finite
loop too shallow for injective self-capture---setting the stage for the
diagonal exclusion of Theorem~\ref{thm:summit}.

\subsection{The \texttt{RepresentationalSystem} Structure}
\label{sec:repsys}

\begin{definition}[\leanref{Basic.lean}: \leanref{RepresentationalSystem}]
\leavevmode\par\noindent
A \emph{representational system} is a tuple $(C, A, \mathrm{represent}, \mathrm{iter\_injective})$ where:
\begin{itemize}
\item $C$ is a category,
\item $A$ is an object of $C$,
\item $\mathrm{represent} : A \to A$ is an endomorphism,
\item $\mathrm{iter\_injective}$: for all $n \neq m$, $\mathrm{represent}^n \neq \mathrm{represent}^m$ in $\mathrm{End}(A)$.
\end{itemize}
\end{definition}

\subsection{Non-Termination and Unboundedness}

\begin{theorem}[\known\ \leanref{Regress.lean}]
\label{thm:regress}
For any representational system $R$:
\begin{enumerate}[label=(\alph*)]
\item (\leanref{regress\_no\_termination}) Distinct indices $n \neq m$ yield distinct meta-level arrows.
\item (\leanref{regress\_iterates\_unbounded}) For every bound $B$, there exists a level $L > B$ with a distinct iterate.
\item (\leanref{regress\_is\_infinite}, \leanref{meta\_range\_infinite}) The set of meta-levels is infinite.
\end{enumerate}
\end{theorem}

\begin{remark}[What if \texttt{iter\_injective} fails?]
\label{rem:dichotomy}
The hypothesis $\mathrm{iter\_injective}$ is \emph{bundled}, not derivable from
a bare endomorphism. An adversary might weaken it: suppose
$\mathrm{represent}^n = \mathrm{represent}^m$ for some $n \neq m$.
\leanref{NoEscapeClosure} formalizes the consequence: let
$e := \mathrm{End.of}(\mathrm{represent})$; if $e^{n+d}=e^n$ with $d>0$ then
every power $e^k$ equals some $e^i$ with $i < n+d$
(\leanref{IterateClosure.End.exists\_lt\_of\_pow\_collision}).
\leanref{IterateClosure.End.pow\_iterate\_dichotomy}
records the \textbf{dichotomy}: for every monoid element~$e$, either
(i)~distinct natural-number powers are all distinct (an infinite injective
$\mathbb{N}$-tower), or (ii)~the set of \emph{distinct} powers is bounded by a
finite index.  In case~(ii) the system has only finitely many representational
layers---it loops.  Either route blocks ``just enough'' self-representation: one
gives an infinite regress; the other gives a finite loop too impoverished to
host injective iterate depth.
\end{remark}

\subsection{The Lawvere Wall}
\label{sec:lawvere}

\begin{theorem}[\known\ \leanref{Lawvere.lean}]
\label{thm:lawvere}
(\leanref{lawvere\_fixed\_point\_Type}) If $s : A \to (A \to B)$ is surjective, then every $f : B \to B$ has a fixed point.
\end{theorem}

\begin{theorem}[\known\ \leanref{LawvereCCCType.lean}]
\label{thm:lawvere-ccc}
(\leanref{lawvere\_fixed\_point\_MonoidalClosedType}) The same conclusion holds when the universal-enumerator hypothesis is recast as surjectivity of $\mathrm{curry}(\mathrm{lawvereBinary}\ s)$ in $\mathrm{MonoidalClosed}(\mathrm{Type}\ u)$.
\end{theorem}

\subsection{The \texttt{OntologicalSlot} Construction}
\label{sec:ontoslot}

\begin{definition}[\novel\ \leanref{Basic.lean}: \leanref{OntologicalSlot}]
\label{def:ontoslot}
For a category $C$ with terminal object and an object $A$:
\[
\mathrm{OntologicalSlot}(C, A) \;:=\;
  \mathrm{endo}(f : A \to A) \;\mid\; \mathrm{obj}(g : \top_C \to A)
\]
an inductive type with two constructors, disjoint by definition.
\end{definition}

\begin{theorem}[\novel\ \leanref{FixedPoints.lean}]
\label{thm:wall}
(\leanref{lawvere\_wall\_is\_not\_dissolution}) For any representational system $R$ and any global element $\mathrm{fp} : \top_C \to A$,
\[
  \mathrm{OntologicalSlot.endo}(R.\mathrm{represent}) \;\neq\; \mathrm{OntologicalSlot.obj}(\mathrm{fp}).
\]
\end{theorem}

\subsection{The Core Master Theorem}

\begin{theorem}[\leanref{Main.lean}]
\label{thm:master}
(\leanref{representational\_regress\_master}) For any representational system $R$:
\[
  \mathrm{representational\_regress\_master\_claim}(R)
\]
i.e., the conjunction of unbounded regress, Lawvere wall non-dissolution, and $T_2$ invariance under homeomorphism.
\end{theorem}

% ============================================================================
% 3. TOPOLOGICAL MODELS
% ============================================================================
\section{Topological Models: Minimal-Case Witnesses}
\label{sec:models}

The general obstruction (Theorem~\ref{thm:generic}) applies to
any pair of spaces with incompatible
\leanref{ChartableR2}-detected boundaries.  This section
establishes the closed cylinder and M\"obius strip as the
\textbf{simplest fully discharged witness pair}: two compact
surfaces whose boundary subtypes are topologically incompatible.

\subsection{The Closed Cylinder}
\label{sec:cylinder}

\begin{definition}[\leanref{CylinderBoundary.lean}]
$\mathrm{ClosedCylinder} := S^1 \times [0,1]$ with the product topology and manifold structure $\mathfrak{R}_1 \times \mathfrak{R}_1^\partial$.
\end{definition}

\begin{theorem}[\known\ \leanref{CylinderBoundary.lean}]
\label{thm:cyl-boundary}
The manifold boundary of \emph{ClosedCylinder} decomposes as:
\[
  \partial(\mathrm{ClosedCylinder}) = F_{\mathrm{bot}} \sqcup F_{\mathrm{top}}
\]
where $F_{\mathrm{bot}} = S^1 \times \{0\}$ and $F_{\mathrm{top}} = S^1 \times \{1\}$ are each connected, and their union is \emph{disconnected}.
\end{theorem}

\subsection{The M\"obius Strip}
\label{sec:mobius}

\begin{definition}[\leanref{CylinderMobius.lean}]
\label{def:mobius-rel}
The \emph{M\"obius equivalence} $\sim$ on $[0,1]^2$ is the equivalence relation generated by
$(0, t) \sim (1, 1-t)$ for $t \in [0,1]$, extended by reflexivity, symmetry, and transitivity.
The \emph{M\"obius strip} is $\mathrm{MobiusStrip} := [0,1]^2 / {\sim}$.
\end{definition}

\begin{theorem}[\leanref{CylinderMobius.lean}]
\label{thm:mobius-props}
\leavevmode
\begin{enumerate}[label=(\alph*)]
\item (\leanref{mobiusRel0Graph\_isClosed}) The graph of $\sim$ is closed in $[0,1]^2 \times [0,1]^2$.
\item (\leanref{CompactSpace MobiusStrip}) $\mathrm{MobiusStrip}$ is compact (quotient of compact).
\item (\leanref{instT2SpaceMobiusStrip}) $\mathrm{MobiusStrip}$ is Hausdorff (closed graph + compact domain).
\item (\leanref{isConnected\_mobiusStripBoundary}) The boundary band $\partial_{\mathrm{model}} := \pi(\{t = 0\} \cup \{t = 1\})$ is connected (one component).
\end{enumerate}
\end{theorem}

\subsection{Boundary Subspace Contrast}

\begin{theorem}[\leanref{MobiusCylinderBoundaryContrast.lean}]
\label{thm:boundary-contrast}
\AfterTheoremHead
\TheoremLeanTag{mobiusStripBoundary\_not\_homeomorphic\_closedCylinderBoundaryUnion}
\[
  \mathrm{IsEmpty}\!\left(\mathrm{Subtype}(\partial_{\mathrm{Mob}}) \simeq_t \mathrm{Subtype}(\partial_{\mathrm{Cyl}})\right).
\]
The M\"obius boundary band (connected) and the cylinder boundary union (disconnected) are not homeomorphic as subspaces.
\end{theorem}

% ============================================================================
% 4. THE CHARTABLER2 FRAMEWORK
% ============================================================================
\section{The \texttt{ChartableR2} Framework}
\label{sec:chartable}

The previous sections established diagonal and categorical barriers.
This section introduces the topological layer: a boundary sensor based on
intrinsic $\mathbb{R}^2$-chartability, together with the general
obstruction lemma (Theorem~\ref{thm:generic}) that will deliver the
class-level topological barrier.

\subsection{Definition and Basic Properties}

\begin{definition}[\leanref{ChartableR2Neighbor.lean}]
\label{def:chartable}
A point $x$ of a topological space $X$ is \emph{$\mathrm{ChartableR2}$} if there exists an open set $U \ni x$ such that $U \cong_t \mathbb{R}^2$.
\end{definition}

\begin{theorem}[\leanref{ChartableR2Neighbor.lean}]
\label{thm:chartable-props}
\leavevmode
\begin{enumerate}[label=(\alph*)]
\item (\leanref{isOpen\_setOf\_chartableR2}) The set of \texttt{ChartableR2} points is open.
\item (\leanref{isClosed\_setOf\_not\_chartableR2}) Its complement is closed.
\item (\leanref{not\_chartableR2\_of\_tendsto\_seq\_atTop}) Sequential stability: limits of non-chartable points are non-chartable.
\end{enumerate}
\end{theorem}

\subsection{Invariance Under Homeomorphism}

\begin{theorem}[\leanref{ChartableR2Bridge.lean}]
\label{thm:chartable-invar}
\AfterTheoremHead
\TheoremLeanTag{Homeomorph.chartableR2\_iff}
If $h : X \simeq_t Y$ is a homeomorphism, then $\mathrm{ChartableR2}(x) \iff \mathrm{ChartableR2}(h(x))$.
\end{theorem}

\subsection{The Generic Boundary-Contrast Obstruction}
\label{sec:generic}

\begin{theorem}[\novel\ \leanref{ChartableR2Bridge.lean}]
\label{thm:generic}
\AfterTheoremHead
\TheoremLeanTag{not\_homeomorphic\_of\_chartableR2\_boundary\_contrast}
Let $X$ and $Y$ be topological spaces with predicates $b_X \subseteq X$ and $b_Y \subseteq Y$ such that:
\begin{enumerate}[label=(\roman*)]
\item $\forall x,\; x \in b_X \iff \neg\,\mathrm{ChartableR2}(x)$,
\item $\forall y,\; y \in b_Y \iff \neg\,\mathrm{ChartableR2}(y)$,
\item $\mathrm{IsEmpty}(\mathrm{Subtype}(b_X) \simeq_t \mathrm{Subtype}(b_Y))$.
\end{enumerate}
Then $\mathrm{IsEmpty}(X \simeq_t Y)$.
\end{theorem}

\begin{proof}[Proof sketch]
Any homeomorphism $e : X \simeq_t Y$ preserves $\mathrm{ChartableR2}$
(Theorem~\ref{thm:chartable-invar}), hence maps $b_X$ to $b_Y$ pointwise.
By $e.\mathrm{subtype}$, we get a homeomorphism
$\mathrm{Subtype}(b_X) \simeq_t \mathrm{Subtype}(b_Y)$, contradicting~(iii).
\end{proof}

Theorem~\ref{thm:generic} subsumes the non-homeomorphism strategy of \S\ref{sec:mfinal}:
any other pair of spaces satisfying the same three hypotheses is also not homeomorphic.

\subsection{Half-Space Obstructions}
\label{sec:halfspace}

\begin{theorem}[\leanref{HalfPlaneVsPlane.lean}]
\label{thm:halfplane}
\AfterTheoremHead
\TheoremLeanTag{euclideanHalfSpace2\_not\_homeomorphic\_euclidean2}
$\mathrm{IsEmpty}(H^2 \simeq_t \mathbb{R}^2)$.
Proof via punctured half-plane (star-convex $\Rightarrow$ simply connected) vs.\ punctured plane (not simply connected, from $\exp$ covering).
\end{theorem}

\begin{theorem}[\leanref{HalfSpaceNeighborVsPlane.lean}]
\label{thm:halfspace-local}
\AfterTheoremHead
\TheoremLeanTag{isEmpty\_homeomorph\_subtype\_openH2\_zero\_nbhd\_subtype\_openPlane}
For open $V \subseteq H^2$ with $0 \in V$ and open $O \subseteq \mathbb{R}^2$, there is no homeomorphism $V \cong_t O$.
\end{theorem}

\subsection{Cylinder C4: Boundary Biconditional}

\begin{theorem}[\leanref{CylinderChartableBoundary.lean}]
\label{thm:cylinder-c4}
\AfterTheoremHead
\TheoremLeanTag{closedCylinder\_boundaryUnion\_iff\_not\_chartableR2}
For all $x \in \mathrm{ClosedCylinder}$:
\[
  x \in \partial_{\mathrm{Cyl}} \;\iff\; \neg\,\mathrm{ChartableR2}(x).
\]
\end{theorem}

% ============================================================================
% 5. THE SEAM-CHART PROBLEM
% ============================================================================
\section{Interior Charts on the M\"obius Witness: The Seam Problem}
\label{sec:seam}

\subsection{Why the Seam Is Hard}
\label{sec:seam-why}

Interior-height points of the fundamental domain fall into three regimes: strict
interior $0<x<1$, where the quotient is injective on small boxes; the vertical
seam $x\in\{0,1\}$ with $t\neq\tfrac{1}{2}$, where trigonometric coordinates can
separate fibers when the height window avoids the equator; and the seam at
$t=\tfrac{1}{2}$, where $(0,\tfrac{1}{2})\sim(1,\tfrac{1}{2})$ and a trigonometric
chart degenerates. $\mathrm{ChartableR2}$ at the equator quotient class is
obtained by a piecewise Euclidean unfolding map (see below).

\subsection{Open-Cell Charts}

\begin{theorem}[\leanref{MobiusChartableBoundary.lean}]
\label{thm:open-cell}
\AfterTheoremHead
\TheoremLeanTag{chartableR2\_mobiusQuotientMk\_of\_Ioo\_square}
If $p \in (0,1)^2$ (strict interior of the fundamental domain), then $\mathrm{ChartableR2}(\llbracket p \rrbracket)$.
\end{theorem}

\subsection{Seam Trig Charts (Away from Equator)}

\begin{definition}[\leanref{MobiusSeamChart.lean}, \leanref{MobiusSeamChartable.lean}]
The \emph{saturated seam patch} $S(t_0, \varepsilon, \delta)$ is the union of the left and right $\varepsilon$-thickened seam strips at height window $(t_0 - \delta, t_0 + \delta)$, saturated under $\sim$.
\end{definition}

\begin{theorem}[\leanref{MobiusSeamTrigInject.lean}]
\label{thm:trig-inject}
\AfterTheoremHead
\TheoremLeanTag{onPatch\_mobiusFDTrigCoord\_eq\_iff\_mobiusRel0}
When $\delta < |t_0 - \tfrac{1}{2}|$, the trigonometric FD coordinate map $\mathrm{mobiusFDTrigCoord}$ satisfies:\newline
\noindent$\mathrm{mobiusFDTrigCoord}(p) = \mathrm{mobiusFDTrigCoord}(q)$ on the saturated patch if and only if $p \sim q$.
\end{theorem}

\begin{theorem}[\leanref{MobiusSeamChartableR2.lean}]
\label{thm:seam-trig}
\AfterTheoremHead
\TheoremLeanTag{chartableR2\_mobiusQuotientMk\_of\_mem\_mobiusSeamSaturatedPatch}
For $t_0 \neq \tfrac{1}{2}$ and appropriate $\varepsilon, \delta$, every point\newline
\noindent in $\pi(S(t_0, \varepsilon, \delta))$ is $\mathrm{ChartableR2}$.
\end{theorem}

\subsection{The Equator Chart: Piecewise Unfolding}
\label{sec:equator}

\begin{definition}[\leanref{MobiusSeamChartableR2.lean}]
\label{def:unfold}
\leavevmode\par\noindent
The \emph{fundamental-domain unfolding map}\newline
\noindent$\mathrm{unfoldFD} : [0,1]^2 \to \mathbb{R}^2$ is:
\[
  \mathrm{unfoldFD}(x, t) :=
  \begin{cases}
    \mathrm{mobiusPlaneCoord}(x, t) & \text{if } x \leq \tfrac{1}{2}, \\
    \begin{aligned}[t]
      &\mathrm{mobiusPlaneTrigSeamPartnerSub}\\
      &\quad\bigl(\mathrm{mobiusPlaneCoord}(x, t)\bigr)
    \end{aligned} & \text{if } x > \tfrac{1}{2}.
  \end{cases}
\]
\end{definition}

\begin{lemma}[\leanref{MobiusSeamChartableR2.lean}]
$\mathrm{unfoldFD}$ is $\sim$-invariant: if $p \sim q$ then $\mathrm{unfoldFD}(p) = \mathrm{unfoldFD}(q)$.
\end{lemma}

\begin{theorem}[\leanref{MobiusSeamChartableR2.lean}]
\label{thm:equator}
\AfterTheoremHead
\TheoremLeanTag{chartableR2\_mobiusQuotientMk\_of\_left\_seam\_half}
$\mathrm{ChartableR2}(\llbracket (0, \tfrac{1}{2}) \rrbracket)$.
\end{theorem}

\subsection{Assembly: All Interior Points Are ChartableR2}

\begin{theorem}[\leanref{MobiusSeamChartableR2.lean}]
\label{thm:interior-height}
\AfterTheoremHead
\TheoremLeanTag{chartableR2\_mobiusQuotientMk\_of\_interior\_height}
For every $p \in [0,1]^2$ with $0 < p_2 < 1$:
\[
  \mathrm{ChartableR2}(\llbracket p \rrbracket).
\]
\end{theorem}

\subsection{M\"obius C4: Boundary Biconditional}

\begin{theorem}[\leanref{MobiusChartableBoundary.lean}]
\label{thm:mobius-c4}
\AfterTheoremHead
\TheoremLeanTag{mobiusStripBoundary\_iff\_not\_chartableR2\_of\_forall\_off\_edge\_chartable}
For all $x \in \mathrm{MobiusStrip}$:
\[
  x \in \partial_{\mathrm{Mob}} \;\iff\; \neg\,\mathrm{ChartableR2}(x).
\]
\end{theorem}

% ============================================================================
% 6. M-FINAL: THE NON-HOMEOMORPHISM
% ============================================================================
\section{M-FINAL: Minimal-Case Witness for the General Obstruction}
\label{sec:mfinal}

The general lemma (Theorem~\ref{thm:generic}) applies to \emph{any} pair of
spaces with incompatible \texttt{ChartableR2}-co-characterized boundaries.
The M\"obius strip and closed cylinder serve as the \textbf{simplest fully
discharged witness}: a concrete compact pair for which all three hypotheses
are proved.  The argument is not specific to these two surfaces; a Klein
bottle vs.\ torus, a projective plane vs.\ a sphere, or any other pair
meeting the same three conditions would yield the same conclusion.

\begin{theorem}[\leanref{MobiusSeamChartableR2.lean}, \leanref{ChartableR2Bridge.lean}]
\label{thm:mfinal}
\AfterTheoremHead
\TheoremLeanTag{mobiusStrip\_not\_homeomorphic\_closedCylinder}
\[
  \mathrm{IsEmpty}\!\left(\mathrm{MobiusStrip} \simeq_t \mathrm{ClosedCylinder}\right).
\]
\end{theorem}

\begin{proof}[Proof sketch]
Instantiate Theorem~\ref{thm:generic} with $b_X = \partial_{\mathrm{Mob}}$,
$b_Y = \partial_{\mathrm{Cyl}}$, the biconditionals
(Theorems~\ref{thm:cylinder-c4} and~\ref{thm:mobius-c4}), and the boundary
incompatibility (Theorem~\ref{thm:boundary-contrast}).
\end{proof}

% ============================================================================
% 7. THE EXTENDED MASTER THEOREM AND INSTANTIATIONS
% ============================================================================
\section{The Extended Master Theorem and Instantiations}
\label{sec:extended}

The summit theorem (Theorem~\ref{thm:summit}) establishes representational
incompleteness as a pure diagonal result.  This section bundles it with the
topological strand and the horn closures into a single verifiable package,
then draws out the consequences for five neighboring debates.

\subsection{Formal Statement}

\begin{theorem}[\leanref{Main.lean}]
\label{thm:extended}
\AfterTheoremHead
\TheoremLeanTag{representational\_regress\_master\_extended}
For any representational system $R$:
\[
  \mathrm{RepresentationalIncompletenessMasterExtended}(R)
\]
where this abbreviation expands to:
\begin{enumerate}[label=(\roman*)]
\item $\mathrm{representational\_regress\_master\_claim}(R)$\newline
\noindent (regress + wall + $T_2$ transport),
\item $\mathrm{IsEmpty}(H^1 \simeq_t \mathbb{R}^1)$ (1D half-space vs.\ line;\newline
\noindent\leanref{representational\_regress\_topology\_halfLineModel}),
\item $\mathrm{IsEmpty}(\mathrm{MobiusStrip} \simeq_t \mathrm{ClosedCylinder})$ (M-FINAL; the minimal-case witness for the general obstruction class).
\end{enumerate}
\end{theorem}

\noindent
This \emph{extends} \leanref{representational\_regress\_master}; it does not replace it.

\subsection{Bundled boundary models}

\begin{definition}[\leanref{ChartableR2Neighbor.lean}]
A \emph{\texttt{ChartableR2} boundary model} on $X$ is a subset $b_X \subseteq X$
together with a proof that
$\forall x,\; x \in b_X \iff \neg\,\mathrm{ChartableR2}(x)$.
\end{definition}

\noindent
\begin{itemize}[leftmargin=1.75em,itemsep=0.35ex,topsep=0.6ex,parsep=0pt]
\item The type \leanref{ChartableR2BoundaryModel} packages such data; the lemma
  \leanref{ChartableR2BoundaryModel.not\_homeomorphic} restates
  Theorem~\ref{thm:generic} in bundled form.
\item The development \leanref{ChartableR2ConcreteBoundaryModels} supplies\newline
  \leanref{mobiusStripChartableR2BoundaryModel} and
  \leanref{closedCylinderChartableR2BoundaryModel}.
\item The same-file theorem\newline
  \leanref{mobiusStrip\_closedCylinder\_boundary\_models\_not\_homeomorphic}\newline
  re-proves M-FINAL along that pipeline.
\end{itemize}

\subsection{Instantiations}
\label{sec:instantiations}

Theorems earlier in the paper are unconditional. This subsection records how the
same statements bear on four debates outside pure mathematics. Where an extra
premise is required, it is stated explicitly; nothing in the formalization
defaults a controversial premise to true.

\subsubsection{Consciousness and self-presence}

\paragraph{Philosophical bridge (premise P2).}
Some accounts distinguish representational self-reference from a stronger notion
of \emph{self-presence} (first-person givenness, non-representational
acquaintance, etc.). Premise~\textbf{P2} is \emph{not} a theorem in the library:
it says only that, for the purposes of a topological comparison, the
representational substrate is modeled by a space whose
\texttt{ChartableR2}-detected boundary has one connectivity pattern (e.g.\
the cylinder's two disjoint components), while the target phenomenon is modeled
by a space whose boundary has a different, incompatible pattern (e.g.\ the
M\"obius strip's single connected band---but the argument is
\emph{not specific to the M\"obius strip}: Theorem~\ref{thm:generic} applies
to \emph{any} pair of spaces with incompatible
\texttt{ChartableR2}-co-characterized boundaries, including Klein bottles,
higher-genus non-orientable surfaces, or any other topological witness that
satisfies the three hypotheses).

\paragraph{What is proved.}
The engine is Theorem~\ref{thm:generic}: \emph{any} two spaces
whose \leanref{ChartableR2}-co-characterized boundaries are
subtype-incompatible are not homeomorphic.  This defines an
open-ended class of obstructed pairs.
Theorem~\ref{thm:mfinal}
($\mathrm{MobiusStrip} \not\simeq_t \mathrm{ClosedCylinder}$)
is the \textbf{minimal-case fully discharged witness}: it
instantiates the general lemma with the simplest compact
example.  Substituting a Klein bottle, projective plane, or
higher-genus non-orientable surface would equally satisfy the
hypotheses; the adversary cannot escape by choosing a different
target.

\paragraph{Closed routes and admissibility.}
Multiple adversarial moves are blocked (see \S\ref{sec:noescape}
for full detail):
\begin{itemize}[leftmargin=1.5em,itemsep=0pt,topsep=0.3ex,parsep=0pt]
\item \emph{Iterate deeper} (Horn~1): the regress is infinite, the Lawvere
  wall persists at every depth---more meta-levels do not close the gap.
\item \emph{Use a finite loop} (Horn~2): the iterate dichotomy shows
  bounded depth cannot host injective self-representation.
\item \emph{Represent via a different object} (Horn~3): section--retraction
  readout collapses to an idempotent on the host and the identity on the
  source; no injective tower on either side.
\item \emph{Use a partial self-model} (Horn~4): the diagonal blind spot
  means the information needed for self-capture is precisely the
  information the partial model cannot represent.
\item \emph{Restrict to decidable predicates} (Horn~6): the $\mathrm{Bool}$
  Lawvere horn is closed; finite cardinality does not help.
\item \emph{Choose a different topological model} (Horn~5): the general
  lemma defines an open-ended class of obstructed pairs---any surface
  with incompatible boundary, not just the M\"obius strip.
\end{itemize}
What remains \textbf{admissible} under P2: the relevant substrate must
\emph{already} have the target topological type (in which case it was
never purely representational in the cylinder sense), or the account
must invoke a discontinuous or non-topological bridge not modeled in
this formalism.  Scaling, iteration, richer self-models, or architectural
reparameterization within a fixed topological type are \emph{inadmissible}
routes to the target.

\subsubsection{Symbol grounding and semantic regress}

In the symbol-grounding problem, arbitrary symbols must be interpreted, and
interpreting the interpretation threatens a regress unless a non-representational
stop is posited. Formally, a representational system already contains the morphism
\texttt{represent} whose iterates index ``one more representational step.'' The
type alias \leanref{SymbolSystem} is definitionally identical to
\leanref{RepresentationalSystem}; lemmas such as
\leanref{symbolGrounding\_metaRegressArrow\_injective} record that distinct
natural-number depths correspond to \emph{distinct} endomorphisms of the
awareness object.

\paragraph{Closed routes and admissibility.}
The horn closures constrain the symbol system on multiple fronts:
\begin{itemize}[leftmargin=1.5em,itemsep=0pt,topsep=0.3ex,parsep=0pt]
\item If iterate indices are injective (Horn~1), interpretation levels
  form an infinite regress with the Lawvere wall intact at every stage.
\item If iterates collide (Horn~2), the system has only finitely many
  distinct interpretation levels---a loop too shallow for full
  self-interpretation.
\item If the system tries a \emph{partial} interpretation function
  (Horn~4), the diagonal exclusion
  (\leanref{lawvere\_diagonal\_not\_in\_range}) guarantees a blind
  spot: the interpretation cannot represent its own diagonal behavior.
\item Restricting to decidable predicates (Horn~6) does not help:
  the $\mathrm{Bool}$-valued Lawvere horn is closed.
\end{itemize}
Admissible grounding must therefore come from \emph{outside} the
representational loop: the system cannot ground itself by iterating
its own interpretive machinery, whether totally or partially, whether
with infinite depth or finite.

\subsubsection{Diagonal arguments and self-interpretation}

Lawvere's theorem in \texttt{Type} implies that there is no family
$s : A \to A \to B$ that realizes every unary $g : A \to B$ as $s(a)$ for some
parameter $a$, once $B$ carries an endomorphism without fixed points.
\leanref{LawvereRegressBridge} collects the instances and connects them to
the categorical presentation;
\leanref{lawvere\_fixed\_point\_stays\_representational} says the
\texttt{OntologicalSlot} wall persists at every iterate depth.

\paragraph{Closed routes and admissibility.}
The diagonal closures are cumulative:
\begin{itemize}[leftmargin=1.5em,itemsep=0pt,topsep=0.3ex,parsep=0pt]
\item \emph{Total self-interpreter} $s : A \to A \to B$: blocked by
  Lawvere for any $B$ with a fixed-point-free endo
  (\leanref{no\_universal\_parametric\_unary\_nat} for $\mathbb{N}$,
  \leanref{no\_universal\_parametric\_unary\_bool} for $\mathrm{Bool}$).
\item \emph{Partial self-interpreter}: abandoning surjectivity does not
  help.  \leanref{lawvere\_diagonal\_not\_in\_range} shows that for
  \emph{any} $s : A \to A \to B$---not necessarily surjective---the
  diagonal $a \mapsto f(s\,a\,a)$ is never $s\,a_0$ for any $a_0$.
  The blind spot is exactly the system's own diagonal behavior.
\item \emph{Restrict to finite / decidable predicates} (Horn~6):
  $\lnot : \mathrm{Bool} \to \mathrm{Bool}$ is fixed-point-free, so the
  two-valued case is no easier.
\item \emph{Climb the tower}: however high the self-applicative coding
  rises, the $\mathrm{endo} \neq \mathrm{obj}$ distinction persists at
  every iterate depth.
\end{itemize}
Admissible self-interpreters must therefore be \emph{incomplete}: they
cannot model their own diagonal behavior.  This is not a resource
limitation (compute, training data, context length) but a structural
one; it holds for every codomain $B$ admitting a fixed-point-free
endomorphism, and for every cardinality of $A$.

\subsubsection{Quantum measurement and observer chains}

Talk of apparatus, environment, and Wigner's-friend hierarchies suggests a chain of
systems, each described relative to a larger ``observer.'' If one ever gives such
a chain a \emph{topological} state space $\mathrm{Obs}$, together with a subset
$\partial_{\mathrm{obs}}$ whose membership agrees with
$\neg\,\mathrm{ChartableR2}$, then Theorem~\ref{thm:generic} applies to any
target space whose boundary subspace is not homeomorphic to
$\partial_{\mathrm{obs}}$---the M\"obius strip is the minimal-case
witness, but the argument works for \emph{any} target space with
incompatible boundary topology (Klein bottle, projective plane, or any
other surface satisfying the three hypotheses of Theorem~\ref{thm:generic}).
\leanref{QuantumObserverChainHypothesis} packages exactly this data as
assumptions; the conclusion
$\mathrm{IsEmpty}(\mathrm{Obs} \simeq_t \mathrm{Target})$ follows for
any target satisfying the hypotheses.

\paragraph{Closed routes and admissibility.}
The same horn closures apply to the observer chain:
\begin{itemize}[leftmargin=1.5em,itemsep=0pt,topsep=0.3ex,parsep=0pt]
\item Iterating the measurement chain (stacking more
  apparatus/environment layers) does not change the topological type of
  the state space (Horns~1--2).
\item Encoding the observer inside a larger system via
  section--retraction does not produce an injective tower (Horn~3).
\item Partial models of the observer--readout process still miss their
  own diagonal (Horn~4).
\item Choosing a different non-orientable target does not evade
  Theorem~\ref{thm:generic} (Horn~5).
\end{itemize}
The chain must either introduce genuinely new topological structure at
some level---not representational iteration of the same type---or
accept that its boundary organization differs from the target.  The
library does \emph{not} define Hilbert spaces, observables, or
projection postulates; any physical application requires an independent
argument that a concrete model satisfies the hypothesis fields.

\subsection{Scaling of language models and similar systems}

Training larger models extends a representational architecture; it does not, by
itself, produce a homeomorphism between a given topological model of that
architecture and a different global model.  Theorem~\ref{thm:generic} is
indifferent to the choice of target: the cylinder/M\"obius pair is the
\emph{minimal-case witness}, but the general lemma blocks every target
space with incompatible \texttt{ChartableR2}-detected boundary.

\paragraph{Closed routes and admissibility.}
\begin{itemize}[leftmargin=1.5em,itemsep=0pt,topsep=0.3ex,parsep=0pt]
\item More parameters, deeper layers, longer context windows: these
  extend the architecture \emph{within} a fixed topological type (Horn~1).
\item Finer-grained self-models built during training: partial
  self-interpreters still miss their own diagonal (Horn~4).
\item Switching from infinite-depth to finite-depth loops (e.g.\
  fixed-iteration transformers): the iterate dichotomy (Horn~2) shows
  bounded depth cannot host injective self-representation.
\item Choosing a different target topology does not help: the general
  lemma covers any pair (Horn~5).
\end{itemize}
Admissible paths to a topological type with different boundary
organization require \emph{structural} change---a different topology
of the underlying state space---not quantitative extension of the same
architecture.

\subsection{Relationship to the Trans-Computational Processing argument}

An earlier line of work argues from computability---halting,
G\"odel-style diagonalization, and related limits---to constraints on
``self-contained'' computational mind models.  Representational
incompleteness (Theorem~\ref{thm:summit}) is logically
\emph{independent}: it requires only a fixed-point-free endomorphism
on the codomain, not recursion-theoretic machinery.  The topological
strand (Theorem~\ref{thm:generic}) is likewise orthogonal: it turns on
boundary-subspace invariants, not on recursiveness.  The two approaches
can reinforce a common \emph{limitative} conclusion about
representation closed under its own rules, but neither formally entails
the other.  Importantly, the present result is \emph{stronger} in one
respect: it holds without any computability assumption on $s$, whereas
computability-based arguments require $s$ to be computable.

\subsection{No-escape analysis: horn-by-horn closure}
\label{sec:noescape}

An adversary confronting the representational regress has, at first glance,
several avenues of retreat.  Each is formalized and shown to fail.  The
numbering below follows the order in which the escape was identified; each
``horn'' names an adversarial move, and the closing theorem blocks it.

\paragraph{Horn~1: Infinite injective tower.}
\emph{Move:} accept $\mathrm{iter\_injective}$ and try to use the tower itself
as a self-representation resource.
\emph{Closure:} Theorem~\ref{thm:regress} yields infinitely many distinct
meta-levels;
Theorem~\ref{thm:wall} (the Lawvere wall) ensures the endo/obj distinction
persists at every level.  The tower does not collapse into closure.

\paragraph{Horn~2: Finite-order / periodic iterates.}
\emph{Move:} weaken $\mathrm{iter\_injective}$---let
$\mathrm{represent}^n = \mathrm{represent}^m$ for some $n \neq m$, so the
loop has only finitely many distinct powers.\newline
\noindent\emph{Closure:}
\leanref{IterateClosure.End.pow\_iterate\_dichotomy}
proves the dichotomy: either the iterates are all distinct (Horn~1) or
there exists a bound $B$ such that every power equals some $e^i$ with
$i < B$.  In the finite case the system's representational depth is
bounded; it cannot host an injective tower at all.

\paragraph{Horn~3: Cross-object readout via section--retraction.}
\emph{Move:} represent $A$ \emph{inside} a different object $B$ using a
section $s : A \to B$ and retraction $r : B \to A$ with $s \circ r = \mathrm{id}_A$.
\emph{Closure:}
\leanref{CrossObjectRepresentation.section\_retraction\_no\_injective\_tower\_either\_side}
shows that $r \circ s$ is idempotent on $B$ (so $(r \circ s)^1 = (r \circ s)^2$,
blocking injective powers on the host side) and $s \circ r = \mathrm{id}_A$
(so all powers are the identity on the source side---trivially periodic).
Neither side supports an injective $\mathbb{N}$-tower.

\paragraph{Horn~4: Partial (non-surjective) self-model.}
\emph{Move:} concede that no \emph{total} self-interpreter exists, but claim
a useful \emph{partial} one.
\emph{Closure:}
\leanref{lawvere\_diagonal\_not\_in\_range}
requires no surjectivity of $s$: for \emph{any} $s : A \to A \to B$ and any
fixed-point-free $f : B \to B$, the diagonal function $a \mapsto f(s\,a\,a)$
is never in the range of $s$.  Every partial self-model necessarily has a
blind spot that includes its own diagonal---the information that matters most
for self-capture is exactly the information it cannot represent.

\paragraph{Horn~5: Different topological models.}
\emph{Move:} reject the cylinder/M\"obius witness; substitute a
Klein bottle, projective plane, or another geometric model.

\emph{Closure:} Theorem~\ref{thm:generic}\newline
(\leanref{not\_homeomorphic\_of\_chartableR2\_boundary\_contrast})\newline
\noindent is \emph{model-agnostic}: it defines an
\textbf{open-ended class} of obstructed pairs---any two spaces
whose \leanref{ChartableR2}-co-characterized boundaries are
subtype-incompatible.  The cylinder/M\"obius pair is the
\emph{minimal-case witness}, not the argument itself.  Replacing
one non-orientable surface with another does not evade the
lemma; it merely changes which boundary incompatibility is
invoked.

\paragraph{Horn~6: Computability / decidable self-reference.}
\emph{Move:} restrict to decidable predicates or finite-valued functions,
where ``self-reference'' seems more tractable than Cantorian diagonals.
\emph{Closure:}
\leanref{no\_universal\_parametric\_unary\_bool}
shows that no $s : A \to A \to \mathrm{Bool}$ can realize every
$g : A \to \mathrm{Bool}$, since $\lnot : \mathrm{Bool} \to \mathrm{Bool}$
is fixed-point-free.  The obstruction is not an artifact of infinite
cardinality; it holds in the decidable, two-valued setting.

\paragraph{Residual openings (not closed here).}
Two further routes are not formally addressed in this paper:
\emph{(i)}~\emph{higher-categorical self-maps} (a self-representation
living in a $2$-category or $\infty$-category, where the monoid of
endomorphisms is replaced by a richer hom-object); and
\emph{(ii)}~\emph{non-categorical frameworks} (a system that resists
categorical encoding entirely).  The first is subject to analogous
monoid-element arguments once a truncation to the $1$-categorical
endomorphism monoid is performed; the second falls outside the scope of
any categorical formalization by definition.  Both are noted as
limitations in \S\ref{sec:limits}.

\subsection{Admissibility: what the closures tell neighboring debates}
\label{sec:admissibility}

Given that Horns~1--6 are closed, the formal landscape constrains what
\emph{admissible} solutions look like in each neighboring debate.  The
reasoning in every case is the same: the closures establish a structural
barrier, so any proposed solution must either (a)~operate outside the
barrier's hypotheses, or (b)~accept the consequences.

\paragraph{Consciousness and self-presence.}
If a system's state-space has one boundary-connectivity pattern (e.g.\
two disjoint components, as in the cylinder), then no homeomorphism maps
it to \emph{any} space with an incompatible pattern (e.g.\ one connected
band, as in the M\"obius strip---but equally a Klein bottle, projective
plane, or any other witness satisfying Theorem~\ref{thm:generic}).  An
admissible account of self-presence must therefore either
(i)~posit that the relevant substrate \emph{already} has the target
topological type, or
(ii)~invoke a discontinuous or non-topological bridge not modeled here.
Scaling---more parameters, longer context windows, deeper
architectures---extends a representational system \emph{within} a fixed
topological type; it does not change boundary organization.

\paragraph{Symbol grounding.}
Any symbol system that iterates interpretation is caught in the
Horns~1/2 dichotomy: either infinite distinct levels (full regress plus
Lawvere wall) or finite-bounded depth (finite loop, no injective tower).
Horn~4 adds that even a \emph{partial} interpretation function misses its
own diagonal.  Admissible grounding must therefore come from
\emph{outside} the representational loop: the system cannot ground itself
by iterating its own interpretive machinery, whether totally or partially.

\paragraph{Diagonal arguments / AI self-interpretation.}
Any proposed self-interpreter $s : A \to A \to B$ (total or partial)
necessarily fails to include its own diagonal
(\leanref{lawvere\_diagonal\_not\_in\_range}); this holds for
$B = \mathbb{N}$ (numeric outputs), $B = \mathrm{Bool}$ (decidable
predicates), and any $B$ with a fixed-point-free endomorphism.
Admissible self-interpreters \emph{must} be incomplete in a precise
sense: they cannot model their own diagonal behavior.  This is not a
resource limitation but a structural one; no amount of compute, training
data, or architectural ingenuity changes the diagonal exclusion.

\paragraph{Quantum measurement chains.}
If the measurement chain's state space has a boundary pattern matching the
cylinder side of the development, then no homeomorphism identifies it with
\emph{any} target whose boundary pattern is incompatible---whether that
target is M\"obius-type, Klein-bottle-type, or any other
space satisfying the generic lemma.  The chain must either introduce
genuinely new topological structure at some level---not representational
iteration of the same type---or accept that its boundary organization
differs from the target.

\paragraph{LLM and AI scaling.}
Scaling extends a representational architecture within a fixed topological
type.  If that type has one boundary-connectivity pattern, scaling does not
produce a homeomorphism to \emph{any} target with an incompatible pattern.
The cylinder/M\"obius pair is the minimal witness; the obstruction is the
general lemma, not the particular pair.
Admissible paths to a different topological type require
\emph{structural} change (different topology of the underlying state
space), not quantitative extension (more of the same).

% ============================================================================
% 8. LIMITATIONS AND SCOPE
% ============================================================================
\section{Limitations and scope}
\label{sec:limits}

\begin{enumerate}
\item \textbf{Hypothesis, not derivation.} $\mathrm{iter\_injective}$ is bundled,
  not proved from minimal data. Not every $f : A \to A$ in every category satisfies it.
  The regress is conditional on this non-triviality hypothesis.

\item \textbf{Lawvere in \texttt{Type}, not abstract CCC.} The formalization uses
  $\mathrm{Type}\ u$ and\newline
  $\mathrm{MonoidalClosed}(\mathrm{Type}\ u)$, not a general closed category over arbitrary~$C$.
  A fully abstract CCC version remains future work.

\item \textbf{\texttt{OntologicalSlot} is definitional, not metaphysical.} The disjointness $\mathrm{endo} \neq \mathrm{obj}$ is trivially true by constructor discrimination. This is the point, but the triviality means the result is \emph{semantic} (about the structure of type theory) rather than a deep theorem.

\item \textbf{Models are chosen, not unique.} The closed cylinder and M\"obius strip are the \emph{minimal-case witnesses} chosen for clarity and tractability.  Theorem~\ref{thm:generic} defines an open-ended class of obstructed pairs; any two spaces with incompatible \texttt{ChartableR2}-co-characterized boundaries (Klein bottle vs.\ torus, projective plane vs.\ sphere, etc.) would serve equally.  Discharging additional witnesses requires proving their boundary biconditionals and subtype incompatibility.

\item \textbf{P2 is philosophical.} The premise ``consciousness requires non-orientable topology'' is not a theorem. Section~\ref{sec:instantiations} states only what follows when P2 is adjoined to the proved non-homeomorphism; accepting or rejecting P2 remains an independent issue.

\item \textbf{No $\mathrm{IsManifold}$ polish.} The M\"obius strip is not given a full Mathlib $\mathrm{IsManifold}$ instance. The proof uses $\mathrm{ChartableR2}$ (open neighborhoods homeomorphic to $\mathbb{R}^2$) as a self-contained predicate, sidestepping the need for a full manifold-with-boundary atlas. This is sufficient for the non-homeomorphism but less elegant than a full atlas approach.

\item \textbf{Route W is incomplete.} An alternative proof via circle-winding / doubling obstructions (Steps~1--4 in \leanref{CylinderMobiusNonhomeo.lean}) is partially formalized. Step~5 (transport under a hypothetical global homeomorphism) is not proved. The flagship proof via $\mathrm{ChartableR2}$ is complete and does not depend on Route~W.

\item \textbf{\texttt{QuantumObserverChainHypothesis}.} This definition only
  bundles topological assumptions (boundary set, \texttt{ChartableR2}
  biconditional,\newline subtype non-homeomorphism to an incompatible target
  boundary---the M\"obius boundary is the minimal witness).  It does not
  import Hilbert spaces, operators, or measurement axioms; a physical
  application would need a separate argument that a concrete model satisfies
  those fields.

\item \textbf{Partial self-models: blind spot, not total impossibility.}
  \leanref{lawvere\_diagonal\_not\_in\_range} proves that every partial
  self-model misses its own diagonal; it does not prove that partial
  self-models are \emph{useless}---only that they are necessarily incomplete
  in a precise, diagonal-shaped way.  Whether that incompleteness matters
  for a given application is a separate question.

\item \textbf{Higher-categorical and non-categorical frameworks.}
  The horn closures rely on $1$-categorical endomorphism monoids.  A
  self-representation living in a $2$-category or $\infty$-category
  replaces the monoid with a richer hom-object; the present arguments
  apply after truncation to the underlying $1$-categorical monoid, but
  a native higher-categorical treatment is not given.  Systems that
  resist categorical encoding entirely fall outside any categorical
  formalization by definition.

\item \textbf{Computability: decidable analogue only.}
  \leanref{no\_universal\_parametric\_unary\_bool} closes the
  $\mathrm{Bool}$-valued Lawvere horn.  A full treatment of Turing-degree
  obstructions (Rice's theorem, Kleene's recursion theorem) in the same
  framework is not attempted; the present result is the finite/decidable
  endpoint of the diagonal family.
\end{enumerate}

% ============================================================================
% 9. RELATED WORK
% ============================================================================
\section{Related work}
\label{sec:related}

\subsection{Formalized topology}
Mathlib's manifold and topology libraries underpin the half-space and chart
lemmas; other proof assistants also host substantial topology.  The weight
here is on a \emph{general} non-homeomorphism lemma (Theorem~\ref{thm:generic})
that defines an open-ended obstruction class, together with a fully discharged
minimal-case witness, rather than on novel classification of differential
invariants.

\subsection{Lawvere's theorem in proof theory}
Lawvere's fixed-point argument appears in multiple formal developments.
\texttt{OntologicalSlot} and the wall lemmas supply complementary
category-internal tagging alongside the \texttt{Type}-level diagonal theorem.

\subsection{Consciousness, G\"odel, and AI}
G\"odel- and computability-based limits on mechanism (including Penrose-style
arguments, discussions of strange loops and self-reference in the sense of
Hofstadter, and information-theoretic theories of
consciousness) occupy a large literature.  Those arguments turn on
decidability or information capacity; ours turns on topological type.
Representational incompleteness (Theorem~\ref{thm:summit}) is
independent of computability: it holds for every codomain admitting a
fixed-point-free endomorphism, regardless of cardinality or
decidability.  The topological strand
(Theorem~\ref{thm:generic}) is likewise orthogonal to G\"odel-style
reasoning: it is a finite-data invariant, not an incompleteness of a
deductive system.

\subsection{Topology in philosophy of mind}
Topological figures of speech appear in some philosophy of mind.  The present
work differs in two ways: (i)~its central topological claims are machine-checked
theorems rather than metaphors, and (ii)~the generalization is explicit---a
general lemma defining an open-ended class of obstructed space-pairs, not a
single hard-wired example.

% ============================================================================
% 10. CONCLUSION
% ============================================================================
\section{Conclusion}
\label{sec:conclusion}

The summit theorem (\leanref{representational\_incompleteness},
Theorem~\ref{thm:summit}) states the central result:
\textbf{no self-modeling system can capture its own diagonal}.  For any
$s : A \to A \to B$ and any fixed-point-free $f : B \to B$, the
diagonal function $a \mapsto f(s\,a\,a)$ is never equal to any row
$s\,a_0$---and if $s$ is claimed to be total and surjective, it simply
cannot exist.  No assumption on the surjectivity, cardinality, or
computability of $s$ is required.  This is
\textbf{representational incompleteness}: the analogue of G\"odel
incompleteness for self-modeling systems rather than formal theories.

Three mutually reinforcing layers make the result robust:

\begin{enumerate}[label=(\roman*)]
\item \textbf{Categorical infrastructure.}\enspace
  A machine-checked regress and Lawvere scaffold, including the
  \leanref{OntologicalSlot} wall that persists at every iterate
  depth, ensures that no climb up the meta-level tower dissolves the
  endo/obj distinction.
\item \textbf{Topological obstruction (general lemma + minimal witness).}\enspace
  Theorem~\ref{thm:generic} defines an \emph{open-ended class} of
  obstructed space-pairs: any two spaces with incompatible
  $\mathrm{ChartableR2}$-co-characterized boundaries are not
  homeomorphic.  The cylinder/M\"obius pair is the fully discharged
  \emph{minimal-case witness}; Klein bottles, projective planes, or
  any other surface satisfying the hypotheses would serve equally.
  The obstacle is \textbf{wrong type}, not
  \textbf{inexhaustible queue}.
\item \textbf{Systematic no-escape closure (Horns~1--6).}\enspace
  Six adversarial retreat routes are formally closed
  (\S\ref{sec:noescape}): infinite regress, finite-order collapse,
  cross-object readout, partial self-models, model substitution, and
  decidable self-reference.  Each closure is a proved theorem in
  Lean~4; together they leave an adversary with no viable route to
  circumvent either the diagonal or the topological barrier.
\end{enumerate}

\paragraph{What is not being claimed.}
The theorems are unconditional mathematics; the \emph{application} to
any particular domain requires domain-specific premises (P2 for
consciousness, \texttt{ChartableR2} biconditionals for a specific
state-space model, etc.).  The formalism does not prove that
consciousness is non-representational, that AI cannot be conscious, or
that symbol grounding is impossible.  It proves that \emph{if} the
relevant state space has the boundary pattern characterized here,
\emph{then} no homeomorphism to a space with a different boundary
pattern exists---and that the six most natural attempts to route around
that barrier fail in formally checkable ways.

\paragraph{The admissibility picture.}
The horn closures (\S\ref{sec:noescape}) together with the
admissibility analysis (\S\ref{sec:admissibility}) leave a clean
landscape.  In each neighboring debate---consciousness,
symbol grounding, diagonal limitations, measurement chains, AI
scaling---the admissible solutions share a common shape:
\emph{structural change} to the topological type of the substrate,
not \emph{quantitative extension} within a fixed type.  Scaling,
iterating, partial modeling, finite-depth loops, cross-object readout,
and decidable restriction are all formally inadmissible.  The
formalism makes this precise and machine-checkable, so that future
claims about circumventing the barrier can be evaluated against a
fixed standard rather than debated informally.

\medskip
\noindent\textbf{Future work.}
\begin{itemize}[leftmargin=1.5em,itemsep=0.15ex,topsep=0.3ex,parsep=0pt]
\item Discharge additional witnesses (Klein bottle, projective plane) for
  the general obstruction class.
\item A full $\mathrm{IsManifold}$ instance for
  $\mathrm{MobiusStrip}$.
\item Completion of the optional Route~W winding track.
\item A more abstract Lawvere statement for general cartesian closed
  categories.
\item Native higher-categorical horn closures ($2$-category and
  $\infty$-category analogues).
\item A full Turing-degree analogue of the decidable Lawvere horn
  (Rice's theorem / Kleene recursion in the same framework).
\end{itemize}

% ============================================================================
% APPENDICES
% ============================================================================
\appendix

\section{Lean Module Map}
\label{app:modules}

\begin{table}[h]
\centering
\footnotesize
\setlength{\tabcolsep}{3pt}
\begin{tabular}{@{}>{\raggedright\arraybackslash}p{3.95cm}>{\raggedright\arraybackslash}p{11.45cm}@{}}
\toprule
\textbf{Module} & \textbf{Role} \\
\midrule
\leanref{Basic} & \leanref{RepresentationalSystem}, iterates, \leanref{OntologicalSlot} \\
\leanref{Regress} & Non-termination, unbounded levels \\
\leanref{FixedPoints} & \leanref{EndoVsPoint}, Lawvere wall non-dissolution \\
\leanref{Lawvere} & Fixed-point theorem in \texttt{Type} \\
\leanref{LawvereCCCType} & \texttt{MonoidalClosed} packaging \\
\leanref{Orientability} & $T_2$ homeomorphism invariance \\
\leanref{CylinderMobius} & Fundamental domain, quotient, compact Hausdorff, boundary band \\
\leanref{CylinderBoundary} & Closed cylinder boundary: two disjoint connected faces \\
\leanref{CylinderChartableBoundary} & Cylinder C4 biconditional \\
\leanref{ChartableR2Neighbor} & \leanref{ChartableR2}, \leanref{ChartableR2BoundaryModel}; openness / stability \\
\leanref{ChartableR2Bridge} & \leanref{Homeomorph.chartableR2\_iff}; \leanref{not\_homeomorphic\_of\_chartableR2\_boundary\_contrast}; \leanref{ChartableR2BoundaryModel.not\_homeomorphic} \\
\leanref{ChartableR2ConcreteBoundaryModels} & M\"obius/cylinder boundary models; conditional observer chain; \leanref{QuantumObserverChainHypothesis} \\
\leanref{SymbolGrounding} & \leanref{SymbolSystem}; \leanref{symbolGrounding\_*} \\
\leanref{LawvereRegressBridge} & \leanref{no\_universal\_parametric\_unary\_nat}, \leanref{no\_universal\_parametric\_unary\_bool}; \leanref{lawvere\_diagonal\_not\_in\_range}; \leanref{lawvere\_fixed\_point\_stays\_representational} \\
\leanref{NoEscapeClosure} & Iterate collisions vs.\ injective towers; section--retraction idempotent readout \\
\leanref{ChartableR2Models} & Interior charts; cylinder interior \leanref{ChartableR2} \\
\leanref{MobiusChartableBoundary} & Open-cell charts; boundary $\neg$\leanref{ChartableR2}; M\"obius C4 packaging \\
\leanref{MobiusSeamChartable} & Saturated seam patches (open, saturated) \\
\leanref{MobiusSeamChart} & IFT-based local model; trig coordinates \\
\leanref{MobiusSeamTrigInject} & Trig injectivity on quotient windows \\
\leanref{MobiusSeamChartableR2} & Seam charts (trig + unfold); interior-height; \textbf{M-FINAL} \\
\leanref{MobiusCylinderBoundaryContrast} & Boundary subspaces not homeomorphic \\
\leanref{CylinderMobiusNonhomeo} & Route W (optional; Steps 1--4) \\
\leanref{HalfLineVsLine} & 1D half-space $\neq$ line \\
\leanref{HalfPlaneVsPlane} & 2D half-space $\neq$ plane \\
\leanref{HalfSpaceNeighborVsPlane} & Local half-space $\neq$ open plane patch \\
\leanref{PuncturedPlaneNotSimplyConnected} & $\mathbb{C} \setminus \{0\}$ not simply connected \\
\leanref{OpenCompactChartObstruction} & Compact $\mathbb{R}^2$ patch $\neq$ $\mathbb{R}^2$ \\
\leanref{Main} & Master theorem + extended bundle \\
\bottomrule
\end{tabular}
\end{table}

\section{Key Theorem Index}
\label{app:theorems}

\begin{table}[h]
\centering
\footnotesize
\setlength{\tabcolsep}{3pt}
\begin{tabular}{@{}>{\raggedright\arraybackslash}p{3.85cm}>{\raggedright\arraybackslash}p{7.55cm}>{\raggedright\arraybackslash}p{3.15cm}@{}}
\toprule
\textbf{Lean name} & \textbf{Statement} & \textbf{Module} \\
\midrule
\leanref{representational\_incompleteness} & \textbf{Summit:} $\forall a_0$, $s\,a_0 \neq \lambda a.\,f(s\,a\,a)$; $\neg\exists$ total $s$ & \leanref{LawvereRegressBridge} \\
\midrule
\leanref{regress\_iterates\_unbounded} & $\forall B, \exists L > B$, distinct iterate & \leanref{Regress} \\
\leanref{lawvere\_fixed\_point\_Type} & Universal enumerator $\Rightarrow$ every endo has fixed point & \leanref{Lawvere} \\
\leanref{lawvere\_wall\_is\_not\_dissolution} & $\mathrm{endo}(\mathrm{represent}) \neq \mathrm{obj}(\mathrm{fp})$ & \leanref{FixedPoints} \\
\leanref{instT2SpaceMobiusStrip} & M\"obius strip is $T_2$ & \leanref{CylinderMobius} \\
\leanref{isConnected\_mobiusStripBoundary} & Boundary band is connected & \leanref{CylinderMobius} \\
\leanref{closedCylinder\_boundary\_eq\_union} & Cylinder boundary = 2 disjoint faces & \leanref{CylinderBoundary} \\
\leanref{euclideanHalfSpace2\_not\_homeomorphic\_euclidean2} & $H^2 \not\cong_t \mathbb{R}^2$ & \leanref{HalfPlaneVsPlane} \\
\leanref{not\_homeomorphic\_of\_chartableR2\_boundary\_contrast} & Generic: incompatible boundaries $\Rightarrow$ $\neg$ homeo & \leanref{ChartableR2Bridge} \\
\leanref{closedCylinder\_boundaryUnion\_iff\_not\_chartableR2} & Cylinder C4 biconditional & \leanref{CylinderChartableBoundary} \\
\leanref{not\_chartableR2\_of\_mem\_mobiusStripBoundary} & M\"obius boundary $\Rightarrow$ $\neg$\leanref{ChartableR2} & \leanref{MobiusChartableBdy} \\
\leanref{chartableR2\_mobiusQuotientMk\_of\_interior\_height} & Interior height $\Rightarrow$ \leanref{ChartableR2} & \leanref{MobiusSeamChartableR2} \\
\leanref{mobiusStripBoundary\_not\_homeomorphic\ldots} & Boundary subtypes not homeo & \leanref{MobiusCylBdyContrast} \\
\leanref{mobiusStrip\_not\_homeomorphic\_closedCylinder} & \textbf{M-FINAL} & \leanref{MobiusSeamChartableR2} \\
\leanref{representational\_regress\_master} & Core master conjunction & \leanref{Main} \\
\leanref{representational\_regress\_master\_extended} & Extended bundle (+ topology) & \leanref{Main} \\
\leanref{QuantumObserverChainHypothesis.not\_homeomorphic\_mobiusStrip} & Conditional: obs.\ chain $\not\simeq_t$ M\"obius & \leanref{ChartableR2ConcreteBoundaryModels} \\
\leanref{no\_universal\_parametric\_unary\_nat} & No universal $s : A \!\to\! A \!\to\! \mathbb{N}$ for all $g : A \!\to\! \mathbb{N}$ & \leanref{LawvereRegressBridge} \\
\leanref{IterateClosure.End.pow\_iterate\_dichotomy} & Injective iterate powers vs.\ finite index bound & \leanref{NoEscapeClosure} \\
\leanref{CrossObjectRepresentation.section\_retraction\_no\_injective\_tower\_either\_side} & Both-side: $\neg$ injective $\mathbb{N}$-powers on source \emph{or} host & \leanref{NoEscapeClosure} \\
\leanref{lawvere\_diagonal\_not\_in\_range} & Diagonal never in range of any $s : A \!\to\! A \!\to\! B$ (no surjectivity) & \leanref{LawvereRegressBridge} \\
\leanref{no\_universal\_parametric\_unary\_bool} & No universal $s : A \!\to\! A \!\to\! \mathrm{Bool}$ for all $g : A \!\to\! \mathrm{Bool}$ & \leanref{LawvereRegressBridge} \\
\bottomrule
\end{tabular}
\end{table}

\section{Reproduction}
\label{app:repro}

\begin{verbatim}
cd representational-incompleteness-lean
lake update
lake exe cache get
lake build RepresentationalIncompleteness
# Verify: 0 sorry
rg "sorry" RepresentationalIncompleteness/
\end{verbatim}

\noindent
Toolchain: \texttt{leanprover/lean4:v4.29.0-rc6}. Mathlib: \texttt{v4.29.0-rc6}.

\end{document}
